% Generated draft: requires mathematical and source review.
\documentclass[11pt]{article}
\usepackage[margin=25mm]{geometry}
\usepackage{fontspec}
\setmainfont{Latin Modern Roman}
\usepackage{amsmath,amssymb,mathtools,unicode-math}
\setmathfont{Latin Modern Math}
\usepackage{enumitem,needspace,xurl,hyperref}
\setlength{\parindent}{0pt}
\setlength{\parskip}{5pt}
\setlength{\emergencystretch}{3em}
\newcommand{\N}{\mathbb{N}}
\newcommand{\Z}{\mathbb{Z}}
\newcommand{\Q}{\mathbb{Q}}
\newcommand{\R}{\mathbb{R}}
\newcommand{\C}{\mathbb{C}}
\newcommand{\F}{\mathbb{F}}
\newcommand{\A}{\mathbb{A}}
\newcommand{\PP}{\mathbb{P}}
\newcommand{\E}{\mathbb{E}}
\newcommand{\eps}{\varepsilon}
\newcommand{\dd}{\,\mathrm{d}}
\begin{document}
\section*{UnsolvedMath ID: 30001943}
\textbf{Source problem code:} catalogue entry
\subsection{Definitions and mathematical statement}
A singularity class here means a class of local types of points, including smooth points, under a specified notion of local equivalence. Write \(\operatorname{Types}(Y)\) for the local types occurring on a space \(Y\), \(\mathcal{NC}\) for the normal-crossings types, and \(X_{\mathrm{nc}}\) for the locus of normal-crossings points of \(X\). The usual local coordinate model for a normal-crossings point is
\[
\{x_1\cdots x_r=0\},\qquad 1\leq r\leq n,
\]
in a smooth ambient space with coordinates \(x_1,\ldots,x_n\); the case \(r=1\) includes smooth points. The ground field, characteristic, permitted spaces, local topology, and treatment of nonsplit branches are not specified in the supplied record, so this model does not settle those choices.

\textbf{Question.} After fixing the intended category and notion of local type, determine a smallest class \(\mathcal S\) containing \(\mathcal{NC}\) such that every variety \(X\) admits a morphism
\[
f\colon X'\longrightarrow X
\]
with
\[
\operatorname{Types}(X')\subseteq\mathcal S,
\qquad
f^{-1}(X_{\mathrm{nc}})\xrightarrow{\;f\;}X_{\mathrm{nc}}
\text{ an isomorphism},
\]
where \(f\) is proper and birational in the algebraic setting, or proper and bimeromorphic in the analytic setting. The quantifier is over every \(X\) in the intended category, and the source \(X'\) and morphism \(f\) may depend on \(X\).

More explicitly, call a class admissible if it contains \(\mathcal{NC}\) and has this modification property for every such \(X\). If ``smallest'' means least under inclusion, the target is an admissible class \(\mathcal S_0\) satisfying
\[
\mathcal S_0\subseteq\mathcal S
\quad\text{for every admissible class }\mathcal S.
\]
The record does not specify whether this least-class interpretation or merely inclusion-minimality is intended, and it does not establish existence of a least class. The formulation therefore remains conditional on these unresolved choices. Although the registry title mentions functorial resolution, the supplied mathematical statement imposes no functoriality requirement.
\subsection{Short English statement}
Determine the smallest collection of allowed local singularity types, containing all normal crossings, that permits every variety to be modified properly and birationally, or bimeromorphically in the analytic setting, while preserving its entire normal-crossings locus. The category, local equivalence convention, and meaning of smallest require clarification. The supplied dated curation labels the problem partially solved; this draft does not independently verify that assessment.
\subsection{Sources}
\begin{itemize}
\item Supplied local UnsolvedMath registry record, permanent ID \(30001943\), source problem code \texttt{catalogue entry}, in the Oberwolfach Reports Open Problems collection. Registry title: \emph{Minimal Singularity Classes for Functorial Resolution}. The proposed year is recorded as 2011; no individual proposer is supplied.
\item Supplied source citation and literature source labeled ``original OWR report'': \emph{Stratified Spaces: Joining Analysis, Topology and Geometry} (2012), Oberwolfach Reports, pp.\ 3217--3286, DOI \texttt{10.4171/owr/2011/56}. \url{https://doi.org/10.4171/owr/2011/56}. The background describes Oberwolfach Reports as documenting mathematical workshops and their open-problem discussions.
\item Supplied literature source labeled ``minimal singularities'': \url{https://doi.org/10.1016/j.aim.2012.08.002}. No authors, full title, or precise supporting theorem are supplied for this entry.
\item Literature assessment dated exactly \texttt{2026-08-21}: the supplied curation labels the problem partially solved and states that the smallest class is determined in low dimensions and low codimension, while the higher-dimensional problem remains open. This is supplied dated curation, not independently verified; the source explicitly requests independent verification of time-sensitive status and cited claims.
\item The registry labels the statement \texttt{corrected\_verified} and reports that its mathematical content was checked against the official Oberwolfach report, with notation, context, and question form normalized. This is a supplied verification claim; no independent source consultation or mathematical verification is asserted here.
\end{itemize}
\end{document}
